\section{Step 4: chopper, first power-on, ESCON setup}\label{sec:walk-chopper}

\subsection{Chopper threshold}
When the motor brakes it pushes energy back into the 48\,V line. The chopper (maxon DSR 70/30) burns it
off above a threshold set by its DIP switches:
\[ V_\mathrm{th} = 75\,\mathrm{V} - \text{(sum of the switches that are ON)} \times 1\,\mathrm{V} \]
with S1 = 1, S2 = 2, S3 = 4, S4 = 8, S5 = 16, S6 = 32.
\begin{enumerate}
\item Mains out. Open the switch cover; use a small flat screwdriver.
\item Set \textbf{S3 and S5 ON}, all others OFF: $75 - (4 + 16) = $ \textbf{55\,V} (Stanford guide p.\,16,
and the three working IBMS towers). ``ON'' is printed on the side opposite the switch numbers.
\end{enumerate}
Why it matters: 55\,V is above the 48\,V supply and below the ESCON's 70\,V maximum. The broken IBMS tower
had all switches OFF (75\,V): it would only act above the ESCON's limit, a plausible cause of that tower's failure.

\textbf{Check:} photograph the switches; S3 and S5 ON.

\subsection{First power-on}
\begin{enumerate}
\item Enable wire (yellow, ESCON J5 pin 2) unplugged so the motor cannot drive. Rope slack.
\item Switch the tower's strip on. Watch for 30\,s: chopper LEDs (shunt active, over-temperature) stay off,
the ESCON LED lights, no smell or noise.
\item Meter, chopper pin 4 to pin 5: 47--49\,V.
\item E-stop (if wired): press, ESCON LED off; release, back on.
\item Strip off; wait a minute before touching wires.
\end{enumerate}
Anything else: strip off and find the cause before trying again (Sect.~\ref{sec:lessons}).

\subsection{ESCON configuration}
Stanford guide ``Motor Driver Setup'' (pp.\,27--32). Done once per driver on the Windows PC.
\begin{enumerate}
\item Install \textbf{ESCON Studio} (maxon website). Close the startup wizard.
\item Enable unplugged, rope slack or off the drum, nobody near the tower.
\item \textbf{USB into ESCON J7 first, then 48\,V on.} The status top right changes to ``connected''
(binoculars button to search if not). Clear any old errors (right-click the error list, ``clear all
entries'').
\item Open \file{escon/moduleA_left_2026-10-01.edc} (our reference: current mode, enable on DigIN2,
setpoint on AnIN2, 0--10\,V = 0--30\,A) and click \textbf{Download} (PC $\rightarrow$ driver).
\item \textbf{Regulation Tuning} $\rightarrow$ auto-tuning, motor free to turn. This adapts the current
controller to this motor.
\item \textbf{Store Parameters.} Without this the driver forgets everything at power-off.
\item Close ESCON Studio (USB stays in). Power-cycle (48\,V off, 30\,s, on), reopen Studio, connect, and check in the Controller Monitor that the settings stayed: current controller, DigIN2 = enable, AnIN2 = set value 0\,V $\rightarrow$ 0\,A, 10\,V $\rightarrow$ 30\,A, rotor position from digital Hall sensors, 30\,A current limit, 4.06\,A nominal, 11 pole pairs. Tuning gains differ from A's (they are per driver). Do not click Download here: it would rewrite them from the PC and hide a failed store.
\item \textbf{Diagnostics}: motor test and Hall sensor test must both pass.
\item Upload (driver $\rightarrow$ PC) and save as \file{escon/module<X>_<side>_<date>.edc}; commit it.
\item 48\,V off first, then USB out. Plug the enable wire back in.
\end{enumerate}
If ESCON Studio offers a firmware update, do not accept it without asking first. If commanding 0\,A later
makes the motor creep or turn easier one way, adjust the offset in ESCON Studio (Stanford p.\,31).

\textbf{Check:} with 48\,V on and the enable low, the ESCON LED \textbf{blinks green slowly} (ready).
Red flashes: count them and look up the code in \file{escon/README.md}.
